% !PRINTED 344
\weberpage{344}
%<<CONT
merkt, eine Annäherung an die absolut kleinste Wurzel. Man kann aber, durch Verlegung des Anfangspunktes, jede Wurzel zur absolut kleinsten machen, wozu freilich die Kenntniss eines bis zu einem gewissen Grade genäherten Werthes nöthig ist.

Eine Formel, die für alle Fälle ausreicht, hat Fr. \gesperrt{Meyer} gegeben\footnote{Mathematische Annalen, Bd. 33.}. Es seien $\alpha_1, \alpha_2, \alpha_3 \ldots \alpha_n$ die Wurzeln der Gleichung $f(x) = 0$, die reell oder imaginär sein können. Man suche einen Punkt $p$ in der $x$-Ebene, so dass sich um $p$ als Mittelpunkt ein Kreis beschreiben lässt, der \gesperrt{nur einen} der Punkte $\alpha_1, \alpha_2 \ldots \alpha_n$ enthält, etwa den Punkt $\alpha_1$, dass also die absoluten Werthe von
\weq{5}{
\frac{p-\alpha_1}{p-\alpha_2}, \quad \frac{p-\alpha_1}{p-\alpha_3}, \quad \cdots \quad \frac{p-\alpha_1}{p-\alpha_n}
}
echte Brüche sind. Solche Punkte existiren immer; sie müssen nöthigenfalls nach der \gesperrt{Sturm}'schen Methode gefunden werden. Wählt man ausserdem noch eine beliebige Function $\varphi(x)$, die mit $f(x)$ keinen gemeinsamen Theiler hat, z.\,B. $\varphi(x) = 1$ und bildet die symmetrischen Functionen
\weq{6}{
y_m = \dfrac{\dfrac{\alpha_1\varphi(\alpha_1)}{(p-\alpha_1)^m} + \dfrac{\alpha_2\varphi(\alpha_2)}{(p-\alpha_2)^m} + \cdots + \dfrac{\alpha_n\varphi(\alpha_n)}{(p-\alpha_n)^m}}{\dfrac{\varphi(\alpha_1)}{(p-\alpha_1)^m} + \dfrac{\varphi(\alpha_2)}{(p-\alpha_2)^m} + \cdots + \dfrac{\varphi(\alpha_n)}{(p-\alpha_n)^m}},
}
so nähern sich diese mit wachsendem $m$ der Grenze $\alpha_1$, und zwar um so schneller, je kleiner die absoluten Werthe der Brüche (5) bereits sind.

Die Richtigkeit hiervon zeigt sich sofort, wenn man die Formel (6) in der Weise schreibt:
\[
y_m = \frac{\alpha_1\varphi(\alpha_1) + \alpha_2\varphi(\alpha_2)\left(\dfrac{p-\alpha_1}{p-\alpha_2}\right)^m + \cdots + \alpha_m\varphi(\alpha_m)\left(\dfrac{p-\alpha_1}{p-\alpha_n}\right)^m}{\varphi(\alpha_1) + \varphi(\alpha_2)\left(\dfrac{p-\alpha_1}{p-\alpha_2}\right)^m + \cdots + \varphi(\alpha_m)\left(\dfrac{p-\alpha_1}{p-\alpha_n}\right)^m}.
\]
% [sic] source prints the final outer subscripts as $m$, although the sequence and the inner $p-\alpha_n$ indicate $n$

\begin{center}
\S.\,111.

\gesperrt{Die Näherungsmethode von Gräffe.}
\end{center}

Auf einem ähnlichen Gedanken, wie die \gesperrt{Bernoulli}'sche Näherung, beruht auch ein Verfahren, das von \gesperrt{Gräffe} an-
%CONT>>
